% !TEX root = ../main.tex
\section{추출 규칙}
\label{app:rules}

\subsection{포스트컨디션}

표~\ref{tab:pc-fields}는 Stacks Blockchain API가 돌려주는 트랜잭션의 포스트컨디션 필드를 \cite{hiroapi} \ref{sec:mapping-rules}절의 규칙에 대응시킨다. 필드 이름은 이 글을 쓰는 시점의 API 스키마를 따랐다. 평가에서는 설정 파일에 API 버전 하나를 고정하고, Originator 모드와 SIP-040의 \emph{may send} 조건이 어떻게 표현되는지 그 버전에서 확인한다.

\begin{table}[htbp]
\centering
\caption{포스트컨디션 필드와 추출에서의 쓰임}
\label{tab:pc-fields}
\small
\begin{tabularx}{\linewidth}{@{}>{\raggedright\arraybackslash}p{0.3\linewidth}>{\raggedright\arraybackslash}X@{}}
\toprule
필드 & 쓰임 \\
\midrule
\code{tx\_status} & 승인은 성공한 트랜잭션에서만 얻는다. 수수료는 실패한 것을 포함해 지갑이 발신한 모든 트랜잭션에 대해 센다. \\
\code{tx\_type}, \code{contract\_call.contract\_id} & 컨트랙트 호출만 쓴다. 여기 적힌 컨트랙트가 진입 컨트랙트 $c(\tau)$이다. \\
\code{sender\_address}, \code{sponsored} & 발신자 $o(\tau)$. 수수료는 대납 여부와 관계없이 발신자의 것으로 센다. \\
\code{fee\_rate} & 수수료 $f(\tau)$. $f_{\max}$에서 자른다. \\
\code{post\_condition\_mode} & Deny와 Originator: 포스트컨디션이 다루는 발신자의 유출만 명시적 승인으로 센다. Allow: 포스트컨디션이 다루지 않는 발신자의 유출을 포괄 승인으로 세고 표시해 둔다. \\
\code{post\_conditions[].principal.type\_id} & \code{principal\_origin}, 또는 발신자 주소와 같은 \code{principal\_standard}: 승인이 될 수 있다. 컨트랙트 principal: 승인이 아니다. \\
\code{post\_conditions[].type} & \code{stx}, \code{fungible}, \code{non\_fungible}. \code{asset}과 함께 자산 $a$를 식별한다. \\
\code{post\_conditions[].condition\_code} & 상한(\code{sent\_equal\_to}, \code{sent\_less\_than}, \code{sent\_less\_than\_or\_equal\_to})과 NFT의 \code{sent} 또는 \emph{may send}: 명시적 승인. 발신자에 대한 하한: 포괄 승인. \code{not\_sent}: 승인 아님. \\
\code{events[]}(자산 전송과 소각) & 실제 유출량 $\mathrm{out}_a(\tau)$: 발신자가 보내거나 소각한 양. \\
\bottomrule
\end{tabularx}
\end{table}

\subsection{Clarity 구문}

표~\ref{tab:clarity-edges}에 분석기가 읽는 구문과 그 구문이 만드는 간선을 나열했다.

\begin{table}[htbp]
\centering
\caption{컨트랙트 $A$의 Clarity 구문과 그 구문이 만드는 간선}
\label{tab:clarity-edges}
\small
\begin{tabularx}{\linewidth}{@{}>{\raggedright\arraybackslash}p{0.42\linewidth}>{\raggedright\arraybackslash}X@{}}
\toprule
$A$ 안의 구문 & 간선 \\
\midrule
\code{(contract-call? 'SP\ldots.b f \ldots)} & 의존 $A\to b$ \\
\code{(contract-call? .b f \ldots)} & 같은 배포자의 $b$로 가는 의존 $A\to b$ \\
\code{(define-constant T 'SP\ldots.b)}와 \code{(contract-call? T f \ldots)} (Clarity~5) & 의존 $A\to b$ \\
\code{(impl-trait 'SP\ldots.d.t)}, \code{(use-trait x 'SP\ldots.d.t)} & trait $t$를 정의한 컨트랙트 $d$로 가는 의존 $A\to d$ \\
\code{as-contract?}, \code{as-contract}, \code{restrict-assets?} 안에서 대상을 알아낸 $b$ 호출 & 위임 $A\to b$ \\
ExecutorDAO 형태인 $A$에서 $b$를 켠 가장 최근의 \code{set-extension} 호출 & 위임 $A\to b$ \\
\code{x}가 trait 타입 인자인 \code{(contract-call? x f \ldots)} & 간선 없음. 동적 디스패치로 기록한다. \\
\code{with-all-assets-unsafe}; 외부 호출을 감싼 \code{as-contract}; \code{as-contract?} 안의 trait 타입 호출 & 위 행들에서 생기는 간선과 별도로 신호로 보여 준다(표~\ref{tab:signals}). \\
\bottomrule
\end{tabularx}
\end{table}

\section{합성 예제}
\label{app:toy}

\ref{sec:method-toy}절의 예제는 난수 시드를 고정한 \code{scripts/toy\_sybil\_example.py}로 만든다. 정직한 컨트랙트 \ToyContracts 개는 SIP-010 trait 하나, 그 trait을 구현하는 토큰 여섯 개, 수학 라이브러리 하나, trait과 라이브러리를 쓰고 \code{as-contract?} 안에서 두 토큰을 내주는 풀 세 개, 세 풀을 호출하는 라우터, 두 토큰을 내주고 라이브러리를 쓰는 볼트, 볼트와 라우터를 확장으로 켠 DAO로 이루어진다. 정직한 지갑 \ToyWallets 개는 저마다 진입 컨트랙트(라우터, 볼트, DAO, 토큰 중) 한 개에서 세 개를 로그정규분포에서 뽑은 값으로 승인하고, $1+\mathrm{Poisson}(8)$개 트랜잭션의 수수료를 낸다. 트랜잭션 하나의 수수료는 최소 수수료에 중앙값이 3인 로그정규분포의 배수를 곱한 값이다. 라우터를 쓰는 지갑은 라우터를 한 번에서 다섯 번 호출하며, 호출마다 풀 하나의 두 토큰을 바인딩한다. 배포 한 번에는 최소 수수료의 \ToyDeployFee 배가 든다. 순위는 공격자의 컨트랙트를 포함해 점수가 대상보다 엄격히 높은 컨트랙트만 세므로, 동점은 대상에게 유리하게 처리되고 표~\ref{tab:toy-cost}의 비용은 실제보다 낮게 나올 수는 있어도 높게 나오지는 않는다. PageRank는 \ref{sec:method}절의 설정을 쓰고, 수수료 상한 $f_{\max}$는 뺐다. 실행 시점 바인딩을 더한 절제 분석에서는 라우터의 바인딩 간선 전체에 라우터의 정적 참조와 같은 가중치를 주고, 그 가중치를 토큰이 바인딩된 횟수에 비례해 나눈다. 표~\ref{tab:toy-ranks}에 공격 크기별 대상의 순위를 적었다.

\begin{table}[htbp]
\centering
\caption{합성 예제에서 공격자 토큰이 전체 컨트랙트 가운데 차지하는 순위(1이 맨 위). $k$는 농장 컨트랙트(A1), 시빌 지갑(A2), 라우터 호출(A4)의 수이다. 열은 표~\ref{tab:toy-cost}의 순위 규칙으로, 차례로 호출자 수, 의존성 PageRank, 균등 재시작으로 Stacks에 옮긴 EndorseRank, \PCER, 실행 시점 trait 바인딩을 더한 \PCER이다.}
\label{tab:toy-ranks}
\small
% scripts/toy_sybil_example.py가 만든 파일이다. 손으로 고치지 않는다.
\begin{tabular}{@{}llrrrrr@{}}
\toprule
공격 & $k$ & 호출자 수 & 의존성 PR & ER(균등) & PC-ER & PC-ER+바인딩 \\
\midrule
A1 & 0 & 10 & 14 & 15 & 15 & 15 \\
A1 & 10 & 10 & 2 & 15 & 15 & 15 \\
A1 & 100 & 10 & 1 & 2 & 15 & 15 \\
A1 & 1,000 & 10 & 1 & 2 & 15 & 15 \\
\addlinespace
A2 & 0 & 10 & 14 & 15 & 15 & 15 \\
A2 & 10 & 10 & 14 & 15 & 15 & 15 \\
A2 & 100 & 3 & 14 & 2 & 15 & 15 \\
A2 & 1,000 & 1 & 14 & 2 & 7 & 8 \\
\addlinespace
A4 & 0 & 10 & 14 & 15 & 15 & 15 \\
A4 & 10 & 10 & 14 & 15 & 15 & 15 \\
A4 & 100 & 10 & 14 & 15 & 15 & 15 \\
A4 & 1,000 & 10 & 14 & 15 & 15 & 9 \\
\bottomrule
\end{tabular}

\end{table}
